% !TeX root = ../main.tex
% Fonte: Lezione 15, pp. 2, 4--5, 9--12. Corretti i passaggi nel rapporto dell'esempio con (2n)!.
\section{Gerarchia degli infiniti}
\label{sec:gerarchia-infiniti}

Per successioni definitivamente positive si scrive $a_n\gg b_n$ quando $a_n/b_n\to+\infty$: $a_n$ cresce più rapidamente di $b_n$ nel senso del rapporto. Equivalentemente $b_n/a_n\to0$. Non si tratta soltanto della disuguaglianza $a_n>b_n$, perché si richiede che il rapporto superi definitivamente ogni soglia positiva.

\begin{proposizione}
Per ogni $q>1$ e $k>0$ fissati,
\[
n^n\gg n!\gg q^n\gg n^k.
\]
In particolare,
\[
\frac{q^n}{n^k}\to+\infty,\qquad
\frac{n!}{q^n}\to+\infty,\qquad
\frac{n^n}{n!}\to+\infty.
\]
\end{proposizione}
\begin{proof}
\textbf{Esponenziale e potenza.} Per $a_n=q^n/n^k>0$,
\[
\frac{a_{n+1}}{a_n}=q\left(\frac n{n+1}\right)^k\longrightarrow q>1.
\]
Il limite della potenza segue dalla continuità in $1$; il criterio del rapporto implica $a_n\to+\infty$. Il ragionamento vale anche per $k$ reale positivo, senza interpretare la potenza come un prodotto di $k$ fattori.

\textbf{Fattoriale ed esponenziale.} Per $b_n=n!/q^n$,
\[
\frac{b_{n+1}}{b_n}=\frac{n+1}{q}\longrightarrow+\infty.
\]
Il criterio del rapporto, nel caso di limite infinito, dà $b_n\to+\infty$.

\textbf{Potenza di esponente variabile e fattoriale.} Per $c_n=n^n/n!$,
\[
\frac{c_{n+1}}{c_n}
=\left(\frac{n+1}{n}\right)^n
=\left(1+\frac1n\right)^n\longrightarrow e>1.
\]
Ancora per il criterio del rapporto, $c_n\to+\infty$.
\end{proof}

La gerarchia si completa con $n^k\gg(\ln n)^b$ per ogni $b>0$, sempre con $k>0$.\footnote{Il confronto fra potenze e logaritmi è enunciato qui; la sua dimostrazione resta rinviata. Non viene usato per dimostrare gli altri confronti o per risolvere gli esercizi successivi.} Il verso dei confronti riguarda il limite per $n\to+\infty$ e dipende dalle ipotesi sulle basi e sugli esponenti.

\begin{esercizio}[Un rapporto con fattoriale e potenza]
Calcolare
\[
\lim_{n\to+\infty}\frac{(2n)!}{2n^n}.
\]
\end{esercizio}
\begin{svolgimento}
Poniamo $a_n=(2n)!/(2n^n)>0$. Si ha esattamente
\[
\begin{aligned}
\frac{a_{n+1}}{a_n}
&=\frac{(2n+2)(2n+1)n^n}{(n+1)^{n+1}}\\
&=2(2n+1)\left(\frac n{n+1}\right)^n
=\frac{2(2n+1)}{(1+1/n)^n}.
\end{aligned}
\]
Il fattore $(n/(n+1))^n$ tende a $1/e$, non a $1$: la base e l'esponente variano insieme. Poiché il denominatore dell'ultima espressione tende a $e>0$, il rapporto tende a $+\infty$. Il criterio del rapporto dà quindi $a_n\to+\infty$.
\end{svolgimento}

\section{Un limite con parametro e termini dominanti}
% Fonte: Lezione 15, pp. 10--12; conservati tutti i casi a<4, a=4, a>4 e il confronto finale con 7.

\begin{esercizio}
Determinare, al variare di $a>0$, il limite
\[
\lim_{n\to+\infty}
\frac{a^n+4^n+n^7+\sqrt{n+1}}
{7^n+n^{53}-\sin(a^{n^2}+n^{11})}.
\]
\end{esercizio}
\begin{svolgimento}
\textbf{Denominatore.} L'oscillazione del seno resta limitata, qualunque sia il comportamento del suo argomento. Raccogliendo $7^n$,
\[
D_n=7^n\left[1+\frac{n^{53}}{7^n}
-\frac{\sin(a^{n^2}+n^{11})}{7^n}\right]
=7^n(1+o(1)).
\]
Il primo rapporto tende a zero per la gerarchia; il secondo è in modulo al più $1/7^n\to0$. La parentesi tende a $1$, quindi è definitivamente positiva e il quoziente è definito su una coda.

\textbf{Numeratore: base maggiore di quattro.} Se $a>4$, la geometrica $(4/a)^n\to0$ e la gerarchia dà $n^7/a^n\to0$. Inoltre $\sqrt{n+1}\le n+1$ per $n\ge1$, quindi anche $\sqrt{n+1}/a^n\to0$. Pertanto
\[
N_n=a^n\left[1+\left(\frac4a\right)^n
+\frac{n^7}{a^n}+\frac{\sqrt{n+1}}{a^n}\right]
=a^n(1+o(1)).
\]
Il rapporto diventa $(a/7)^n$ moltiplicato per un fattore tendente a $1$: tende a zero se $4<a<7$, a $1$ se $a=7$, a $+\infty$ se $a>7$.

\textbf{Numeratore: base minore o uguale a quattro.} Se $0<a<4$, raccogliendo $4^n$ si ottiene $N_n=4^n(1+o(1))$, perché $(a/4)^n\to0$ e i termini di potenza e radice sono trascurabili rispetto a $4^n$. Il quoziente tende quindi a zero come $(4/7)^n$.

Se $a=4$, i due termini esponenziali coincidono e si sommano: $N_n=2\cdot4^n(1+o(1))$. Il quoziente tende ancora a zero, perché $2(4/7)^n\to0$.

Riunendo tutti i casi, il limite è
\[
\boxed{\begin{cases}
0,&0<a<7,\\
1,&a=7,\\
+\infty,&a>7.
\end{cases}}
\]
\end{svolgimento}
